\chapter{Casos Especiais de Citação}
\label{chap:casos-citacao}

As citações devem permitir que o leitor identifique a fonte utilizada e compreenda a relação entre a ideia citada e o argumento do trabalho. Este capítulo apresenta situações frequentes que exigem atenção.

Os nomes, datas, títulos e trechos utilizados nos exemplos deste capítulo são fictícios e têm finalidade exclusivamente didática, salvo quando uma fonte real for expressamente identificada. Autores fictícios usados apenas para demonstrar chamadas não exigem inclusão na lista final de referências.

\section{Um autor}

Quando a obra possui um autor, a chamada deve indicar sobrenome e ano, conforme o sistema adotado. A citação pode aparecer integrada ao texto ou entre parênteses.

Exemplo de chamada integrada: Segundo Silva (2024), a delimitação do problema orienta as escolhas metodológicas.

Exemplo de chamada parentética: A delimitação do problema orienta as escolhas metodológicas (Silva, 2024).

\section{Dois ou três autores}

Quando houver dois ou três autores, todos devem ser indicados na chamada, respeitando a forma adotada pela norma e pelo padrão institucional.

Exemplo: A escrita acadêmica exige planejamento, revisão e domínio das fontes (Silva; Almeida; Costa, 2024).

\section{Mais de três autores}

Em obras com mais de três autores, pode-se indicar o primeiro autor seguido da expressão adequada, conforme a norma adotada. O estudante deve manter o mesmo padrão em todo o trabalho.

Exemplo: A revisão bibliográfica deve seguir critérios explícitos de seleção (Silva et al., 2024).

\section{Autor entidade}

Quando a autoria é de instituição, órgão, associação, universidade ou entidade coletiva, a chamada deve identificar a entidade responsável. Siglas podem ser usadas quando forem conhecidas e estiverem apresentadas adequadamente.

Exemplo: A apresentação de trabalhos acadêmicos deve seguir critérios de estrutura e normalização (Associação Brasileira de Normas Técnicas, 2024).

\section{Fonte sem autoria identificada}

Quando não houver autor identificado, a chamada deve ser feita pelo título ou pela primeira palavra significativa do título, conforme a regra aplicável. Fontes sem autoria devem ser usadas com cuidado e avaliadas quanto à confiabilidade.

\section{Autores com mesmo sobrenome}

Quando autores diferentes possuem o mesmo sobrenome e publicações no mesmo ano ou no mesmo contexto de citação, é necessário acrescentar elementos que permitam diferenciá-los, como iniciais dos prenomes, conforme o sistema adotado.

\section{Mais de uma obra do mesmo autor}

Quando o mesmo autor possui mais de uma obra no mesmo ano, usam-se letras após o ano para diferenciar as obras, desde que a lista de referências também mantenha essa distinção.

Exemplo: Silva (2024a) discute a delimitação do problema, enquanto Silva (2024b) aborda a construção dos objetivos.

\section{Citação de várias obras}

Quando uma afirmação é sustentada por várias obras, as fontes devem ser organizadas conforme o padrão normativo adotado. O estudante deve evitar acúmulo desnecessário de referências e selecionar as mais pertinentes.

Exemplo: A revisão de literatura deve ser seletiva e articulada ao problema (Almeida, 2022; Costa, 2023; Silva, 2024).

\section{Citação direta com supressão}

Quando parte do trecho citado é omitida, a supressão deve ser indicada de forma padronizada. A omissão não pode alterar o sentido original.

Exemplo: ``A pesquisa acadêmica [...] exige coerência entre problema, método e análise'' (Silva, 2024, p. 25).

\section{Citação direta com acréscimo}

Quando o estudante acrescenta palavra ou explicação dentro de citação direta, deve marcar que o acréscimo não pertence ao texto original.

Exemplo: ``A instituição [universidade] deve definir critérios de apresentação'' (Silva, 2024, p. 30).

\section{Citação direta com grifo}

Quando houver destaque em trecho citado, deve-se indicar se o grifo pertence ao original ou se foi acrescentado pelo autor do TCC.

Exemplo: ``O método deve ser compatível com o problema'' (Silva, 2024, p. 41, grifo nosso).

\section{Citação traduzida}

Quando o estudante traduz trecho de obra em outro idioma, deve informar que a tradução é sua, conforme o padrão adotado.

Exemplo: ``A revisão deve ser transparente e reprodutível'' (Smith, 2023, p. 15, tradução nossa).

\section{Informação verbal}

Informações obtidas oralmente, como palestras, entrevistas informais ou comunicações pessoais, devem ser usadas com cautela. Quando forem relevantes, a origem deve ser indicada no texto ou em nota, conforme orientação institucional.

O estudante deve avaliar se a informação verbal pode ser verificada. Quando não puder, talvez seja melhor não usá-la como sustentação central do argumento.

\section{Documento não paginado}

Em documentos digitais sem paginação, a localização pode ser indicada por seção, capítulo, parágrafo, tempo de vídeo ou outra forma que ajude o leitor a localizar o trecho, quando necessário.

\section{Citação de legislação}

Na citação de legislação, deve-se identificar o ente responsável, o tipo de ato, número, data e informações necessárias para localização. No texto, evite citar apenas o número da lei sem contextualizar seu conteúdo.

\section{Citação de dados e bases}

Dados estatísticos, bases públicas e levantamentos institucionais também precisam de fonte. Informe órgão, ano, tabela, base, painel, relatório ou recorte utilizado.

\caixaorientacao{O estudante deve padronizar as citações desde o início. Corrigir dezenas de chamadas no fim do trabalho é mais trabalhoso e aumenta o risco de inconsistência.}
